%! Counting and Probability with Counting — Unit 3, Lesson 3.5 — Student Packet Cover Sheet
\documentclass[12pt]{article}
\usepackage{saar-article}
\usepackage{saar-boxes}
\usepackage{ltablex}
\keepXColumns
% Measured banner: the band grows to fit the title block, so a title long enough to wrap grows
% the band rather than running out of it (the stats course's \coverbanner; saar-article does not
% carry one, so it is defined here — every cover copies this block verbatim).
\newsavebox{\bannerbox}
\newlength{\bannerht}
\newlength{\coverbannerpad}
\setlength{\coverbannerpad}{0.16in}       % breathing room above and below the text
\newcommand{\coverbanner}[2]{%
  \sbox{\bannerbox}{%
    \begin{minipage}{\dimexpr\paperwidth-1.4in\relax}%
      \centering
      {\color{white}\textbf{\LARGE \CourseName}}\\[4pt]
      {\color{frostmid}\large\textbf{#1}}\\[6pt]
      {\color{white}\textbf{\large #2}}%
    \end{minipage}}%
  \setlength{\bannerht}{\dimexpr\ht\bannerbox+\dp\bannerbox+2\coverbannerpad\relax}%
  \noindent\begin{tikzpicture}[remember picture, overlay]
    \fill[cerulean] (current page.north west)
      rectangle ([yshift=-\bannerht]current page.north east);
    \node[anchor=north, inner sep=0pt]
      at ([yshift=-\coverbannerpad]current page.north) {\usebox{\bannerbox}};
  \end{tikzpicture}\par
  % The text body starts 0.6in below the page edge (geometry top=0.6in), so reserve only what
  % the band overruns that, plus a gap under the band.
  \vspace*{\dimexpr\bannerht-0.6in+0.16in\relax}%
}

\begin{document}

% Full-bleed cerulean banner, sized to the title block.
\coverbanner{Unit 3}{Lesson 3.5 \quad Counting and Probability with Counting}

\vspace{0.06in}
% NAMESTRIP: this is the ONLY component that carries the name row. Everything else
% is stapled behind it.
\namedateperiod
\vspace{0.18in}

% GRADUAL RELEASE: the vocabulary is taught in the guided notes, so the targets name it
% plainly. The standard codes stay in the teacher-facing plan.
\begin{learningtargetbox}
\small
\begin{enumerate}[leftmargin=1.5em, itemsep=4pt]
  \item \ldots{} use the \textbf{Fundamental Counting Principle} and \textbf{factorial} notation
        to count the outcomes of several decisions.
  \item \ldots{} calculate a \textbf{permutation} $_nP_r$ when order matters, and verify it
        with a calculator.
  \item \ldots{} calculate a \textbf{combination} $_nC_r$ when order does not matter, and
        explain why dividing by $r!$ removes the repeated orderings.
  \item \ldots{} decide whether a situation is a permutation or a combination by checking
        whether order matters.
  \item \ldots{} find a probability as favorable outcomes over total outcomes using
        combinations, and make a prediction from it.
\end{enumerate}
\end{learningtargetbox}

\vspace{0.12in}

\begin{tocbox}
\small
\renewcommand{\arraystretch}{1.5}
\begin{tabularx}{\linewidth}{@{} c l X @{}}
\rowcolor{frost}
\textbf{\#} & \textbf{Component} & \textbf{Description} \\
\midrule
1 & Warm-Up & Products and fractions, cancelling first \\
2 & Guided Notes \& Practice & Student council at Fairview High: the Counting Principle,
             factorials, permutations, combinations, and offices versus a committee --- then
             a playlist we work together \\
% Row 3 is the lesson's GRADED practice, assigned at the close. The row never changes with the form.
3 & Homework & Locker codes and a fantasy-team draft: the counts, the probability of an all-rookie
             team, and why the question decides the count. \textbf{Scored 1 to 5, assigned
             today, due the first class after two study halls.} \\
\end{tabularx}
\end{tocbox}

\vspace{0.12in}

% Keep in Mind carries the lesson's CONTENT — the rule, the test, the trap — never its process.
\begin{remindbox}
\small
\textbf{Check order before you count.} Swap two people or objects: if the outcome changes, order
matters and the count is a \textbf{permutation}, $_nP_r = \dfrac{n!}{(n-r)!}$; if it does not
change, order does not matter and the count is a \textbf{combination},
$_nC_r = \dfrac{n!}{r!\,(n-r)!}$. Both count without repetition; the \textbf{Fundamental Counting
Principle} also handles choices that may repeat. A probability with counting is
$\dfrac{\text{favorable}}{\text{total}}$, with the same kind of count on top and bottom.
\textbf{Watch out:} the same $n$ and $r$ give different answers when the question changes, and
the word \emph{choose} does not settle it. The roles do.
\end{remindbox}

\end{document}
